% !TEX root = main.tex


\section{Results and perspectives} \label{sec:results}
We evaluate the predictive accuracy of the MJPs under the model settings from \cref{sec:mjp}, compare against the reference predictors in \cref{sec:reference_forecasts}, and analyze differences in transient distributions and transition residuals. In addition to error scores, calibration and discrimination of the forecasts are assessed using reliability diagrams and checkerboard plots, and sharpness through an entropy measure. Further, we propose a use case for the forecasting models, i.e., a decision making framework taking the probabilistic forecasts as inputs, demonstrating how prediction methods can be incorporated in maintenance and inspection planning.

\cref{fig:flowchart} provides a flowchart of this present study, displaying how its different parts  are connected.

\begin{figure}[h!]
\centering
% !TEX root = main.tex


\newcommand\distvert{0.75cm}
\newcommand\distlong{0.5cm}
\newcommand\boxwidth{4cm}
\newcommand\redboxheight{2.25cm}
\newcommand\otherboxheight{1.5cm}
\newcommand\cornersround{1mm}
\usetikzlibrary{fit}


\begin{tikzpicture}[
roundnode/.style={rectangle, draw=black, minimum size=5mm,text width=\boxwidth, minimum height=\otherboxheight, align = center,draw=black!100, fill=black!0},
roundnode4/.style={rectangle, draw=black, minimum size=5mm,text width=1.25*\boxwidth, minimum height=0.5*\otherboxheight, align = center,draw=black!100, fill=black!0},
%rounded corners=\cornersround
roundnode2/.style={rectangle, draw=black, minimum size=5mm,text width=8.5cm, minimum height=\otherboxheight, align = center,draw=black!100, fill=black!0},
%
roundnode3/.style={rectangle, draw=black, minimum size=5mm,text width=.1cm, minimum height=\otherboxheight, align = center,draw=white!100, fill=black!0},
%
squarednode1/.style={rectangle, draw=red!20, fill=red!3, very thick,minimum size=5mm,text width=\boxwidth, minimum height=\redboxheight, align=center, anchor=north,rounded corners=\cornersround},
squarednode2/.style={rectangle, draw=cyan!20, fill=cyan!3, very thick,minimum size=5mm,text width=\boxwidth, minimum height=\otherboxheight, align = center,rounded corners=\cornersround},
invisible/.style={rectangle, draw=white, fill=white, minimum size=0mm}
]

 


\node[roundnode] (data) at (0,0) {Defect data with covariates and censoring (\cref{sec:defect_trajectories}). };
\node[roundnode] (MJP) [right=\distvert of data] {Markov jump process (\cref{sec:mjp}).};

\node[roundnode4, anchor=north west] (generators) at ([xshift=1.5*\distvert]MJP.north east) {Generator parameterizations (\cref{sec:generator_params}).};
\node[roundnode4] (sojourn) [below=0.5*\distvert of generators] {Modified sojourn time estimates (\cref{sec:warping}).};
\node[roundnode4] (regressions) [below=0.5*\distvert of sojourn] {Exogenous dependencies (\cref{sec:regression,sec:regression_state_dependent}).};
\node[roundnode4] (mixture) [below=0.5*\distvert of regressions] {Mixture distributions (\cref{sec:mixture}).};

\node[draw, dotted, thick, rounded corners=\cornersround, fit=(regressions)(generators)(sojourn)(mixture), inner sep=2mm]  (MJP_settings_box) {};
\node[rotate=-90, anchor=south, xshift=2mm] at (MJP_settings_box.east) {\small MJP settings};


\node[roundnode] (estimation) [below=\distvert of MJP] {Optimum score estimation (\cref{sec:optimum_score_estimation}).};
\node[roundnode] (prediction) [below=0.5*\distvert of estimation] {Probabilistic prediction using transient distribution (\cref{sec:transient_dist}).};
\node[draw, dotted, thick, rounded corners=\cornersround, fit=(estimation)(prediction), inner sep=2mm]  (estim_pred_box) {};
\node[rotate=-90, anchor=south, xshift=2mm] at (estim_pred_box.east) {\small Cross validation};

\node[roundnode] (reference) [left=\distvert of estim_pred_box, xshift=2mm] {Reference predictors (\cref{sec:reference_forecasts}).};
    
\node[roundnode3] (hidden) [below=\distvert of estim_pred_box] {};

\node[roundnode] (evaulation) [left=\distvert/2-0.05cm of hidden] {Assessment of predictive performance (\cref{sec:error_metrics,sec:prediction_error,sec:results_transient_dist_markov,sec:reliability_analysis}).};
\node[roundnode] (decisions) [right=\distvert/2-0.05cm of hidden] {Decision making framework (\cref{sec:proposition_for_use}).};


\draw[->] (data) to (MJP);
\draw[->] (MJP) to (estim_pred_box);
\draw[<->] (MJP.east) to (MJP_settings_box);
\draw[<->] (reference) to (estim_pred_box);
\draw[<->] (estimation) to (prediction);
\draw[->] (prediction) to (evaulation);
\draw[->] (prediction) to (decisions);
\draw[->] (reference) to (evaulation);



\end{tikzpicture}

\caption{Workflow of the present study.} \label{fig:flowchart}
\end{figure}


\subsection{Prediction accuracy} \label{sec:prediction_error}
The predictions will be evaluated using the metrics RPS and log score (see \cref{sec:error_metrics}). This section outlines the procedure for how this has been done, including obtained diagnostics from the model estimation. Finally, model performances are compared based on the computed measures of predictive accuracy.

\subsubsection{Out-of-sample evaluation via cross-validation} \label{sec:CV}
The calculation of prediction accuracy is done within a $k$-fold cross-validation (CV) procedure, so that the predictions are evaluated \textit{out-of-sample}. This will be described in the following.

\begin{enumerate}
    \item At random, divide the set of defects into $k$ subsets of approximately equal size given a fixed seed. To avoid leakage (i.e. that defects coming from nearby locations are not in both the training and test set), define a group identifier $g(s) = (l^{(s)},\left \lfloor{p^{(s)}/\Delta}\right \rfloor )$, with $l^{(s)}$ denoting the rail line of defect $s$, $p^{(s)}$ its kilometric position, and $\Delta$ the tolerance threshold for spatial overlap. Let $\mathcal{G} = \{g(s) : s \in \mathcal S\}$ denote the set of unique groups of defects. We then construct a random partition $\{\mathcal{G}_1, \dots, \mathcal{G}_k\}$ (approximately equal in number of groups) and define $\mathcal S_i = \{s : g(s) \in \mathcal G_i\}$. 
    \item For each of the $k$ subsets (folds):
    \begin{enumerate}
        \item Initialize $i = 1$. By $\mathcal{S}_i$, we denote the defects belonging to the $i$'th fold. The set of remaining defects is denoted $\mathcal{S}_{\neg i}$, such that $\mathcal{S} = \mathcal{S}_i \cup \mathcal{S}_{\neg i}$ and $\mathcal{S}_i \cap \mathcal{S}_{\neg i} = \emptyset$.
        \item Model estimations are performed on the set $\mathcal{S}_{\neg i}$. The parameters of the Markov jump process are determined as described in \cref{sec:optimum_score_estimation}, and the parameters of the cumulative link model are determined as described in \cref{sec:cumulative_link_model}.
        \item The models are used to produce predictions of the distributions of $\widehat{\bm{\pi}}^{(s)}\left( t_2^{(s)}\right)$ for the defects $s \in \mathcal{S}_i$ provided their corresponding initial state $y_1^{(s)}$, their initial time stamp $t_1^{(s)}$, and the covariates $\bm{z}^{(s)}$ (for the Markov jump process) or $\bm{z}'^{(s)}$ (for the cumulative link model).
        \item If $i < k$, then increase $i = i + 1$ and GOTO (a); else GOTO 3.
    \end{enumerate}
    \item For each model with the corresponding set of predictions, compute the prediction scores given all end point observations $y_2^{(s)}$ for $s \in \mathcal{S}$ for each metric described in \cref{sec:error_metrics}.
\end{enumerate}
 
\subsubsection{Computational settings and recorded diagnostics} \label{sec:CV_details}
The larger $k$ is chosen in the CV, the higher the computational burden becomes \citep{fushiki2011a}. In this study, $k=5$ is used, which given the procedure in \cref{sec:CV} results in folds containing \num{703}, \num{725}, \num{751}, \num{752}, and \num{727} defects, respectively.\footnote{The seed for the CV is set at \texttt{1}.} The threshold for the spatial overlap is $\Delta = 0.1$ in kilometers.

For the DTMC (cf. \cref{sec:DTMC}), we use $\alpha=0.1$ and discretize the observation interval $\tau$ (in years) using the bin edges
$
0,0.05,0.10,\dots,3.50,\infty
$.
For the regressions in \cref{eq:trans_rates_covariates,eq:trans_rates_covariates2}, the link function $g(x) = \exp(x)$ is used\footnote{Using $g(x)=\log(1+\exp x)$ (softplus) gave similar results as the exponential link. The square link $g(x)=x^2$ was less numerically robust in our experiments, likely because it is not an injective mapping.} across all MJP specifications.
For the mixture MJPs, the selection of $G$, that is, the number of latent components (see \cref{sec:mixture}) is chosen through a simple iterative procedure, starting by $G=2$ and increasing until AIC stops improving. We use $G = 4$ for all generator parameterizations. Remark that a small ridge regularization was imposed on the objective for the mixture MJPs, i.e. $\mathcal D_{|\mathcal S_i|} (\bm \theta) - \lambda_{\text{ridge}} ||\bm \theta||_{2}^2$ with $\lambda_{\text{ridge}} = 0.001$.  This improves numerical robustness by penalizing extreme parameter values, since mixture-model estimation is prone to multiple local optima \citep{shireman2016a}.

Starting values for the parameters in the numerical optimization are as follows: 1) all free transition rates parameters $\lambda_{i,j}^{(g)}$ for $g = 1,\dots,G$ were initialized at $\exp(-1)$ given any generator parameterization; 2) all regression coefficients $ \beta_k$ for $ k = 1,\dots, \widetilde{p}$ were initialized at $0$ where $\widetilde{p}$ denotes the number of covariate coefficients; 3) the sojourn time weights were initialized to $\xi_i=1$ with $\xi_m$ fixed at $1$ (using the working parameters $\widetilde \xi$ in \cref{sec:warping}); and 4) for mixture MJPs, the unconstrained mixture weights were initialized equally (implying $\tilde\kappa_g = 1$ for $g=1,\dots,G-1$ with $\tilde\kappa_G=1$), which corresponds to uniform mixture weights $\kappa_g = 1/G$ under the normalization $\kappa_g=\tilde\kappa_g/\sum_{i=1}^G \tilde\kappa_i $ (\cref{sec:mixture}).  
These starting values were used across folds, so that performance across folds is evaluated wrt. robustness to data subsetting and not optimizer initialization. However, numerical experiments showed that model performances were similar across random initializations.

Across folds, estimation diagnostics are recorded; Appendix \ref{sec:diagnostics_estimation} provides a fold-wise overview of the gradient norms of the objective function in \cref{eq:mean_score_combination} and Akaike information criterion (AIC)\footnote{AIC is calculated as $\text{AIC} = 2 (\text{dim}_{\bm \theta} - \ell (\widehat{\bm \theta}  ; \mathcal{S}_i) )$ for fold $i$. Note that $\widehat{\bm \theta}$ are determined using the objective function in \cref{eq:mean_score_combination}, but AIC is computed from the log-likelihood, $\ell$, evaluated at $\widehat{\bm\theta}$.} given the obtained model parameters for each fold and MJP specification.
Appendix \ref{sec:diagnostics_estimation} shows that the numerical optimization is generally stable across folds and models, with small gradient norms at the obtained solutions. The optimizer’s convergence codes indicate convergence for all models on all folds, except three cases of the mixture MJP, that is: \(\bm{A}_{\RN{1}}\) with covariates on fold 2, \(\bm{A}_{\RN{2}}\) with covariates on fold 5, and \(\bm{A}_{\RN{4}}\) without covariates on fold 1 (also reflected by the gradient norms). For these cases, convergence could be reached by adjusting the initialization; but for consistency, we report results based on the initialization described above. As will appear in \cref{sec:prediction_scores_tables}, the AIC values in Appendix \ref{sec:diagnostics_estimation} are by and large lowest for the models with the highest predictive accuracy.
%Subsequently (and because this study focuses on predictive performance), emphasis will be placed on the predictive scores rather than AIC in the model evaluation.
%%#187       gerlang_cov1_exp_exp_all_no_warp_mixture4    2         
%%#260 gerlang_relax_cov1_exp_exp_all_no_warp_mixture4    5        
%%#296   tridiagonal_cov0_exp_exp_all_no_warp_mixture4    1    

\subsubsection{Summary of predictive scores across models} \label{sec:prediction_scores_tables}
\cref{tab:prediction_scores} reports the mean predictive scores across folds from the cross-validation procedure (\cref{sec:CV,sec:CV_details}), together with their standard deviations across folds.
An ensemble forecast is constructed by averaging all predictive distributions from all other predictors in the table.

\begin{longtable}[t]{>{}lllccc}
\caption{\label{tab:prediction_scores}Prediction scores from $k$-fold cross-prediction procedures in terms of average ranked probability score and average log score, with standard deviations across folds in parentheses. ``Exogenous info.'' indicates whether the rail characteristics described in \Cref{sec:data} are included: Yes (global) uses \cref{eq:trans_rates_covariates}; Yes (state-spec.) uses \cref{eq:trans_rates_covariates2}. ``\# Pars.'' indicates the number of estimated model parameters. Bold numbers indicate best performance within a model class (where applicable).}\\*
\toprule
 & Model & Exogenous info. & \# Pars. & RPS & Log S.\\*
\midrule
\endfirsthead
\caption[]{\textit{(continued)}}\\*
\toprule
 & Model & Exogenous info. & \# Pars. & RPS & Log S.\\*
\midrule
\endhead

\endfoot
\bottomrule
\endlastfoot
\addlinespace[0.3em]
\multicolumn{6}{l}{\textbf{\textbf{Naïve predictor}}}\\*
\hspace{1em} & Uniform distribution & No & -- & 0.6029 (0.0093) & -1.3027 (0.0144)\\*
\cmidrule{1-6}
\addlinespace[0.3em]
\multicolumn{6}{l}{\textbf{\textbf{Cumulative link model}}}\\*
\hspace{1em} & Logistic link & No & 9 & 0.4551 (0.0108) & -1.1919 (0.0173)\\*
\nopagebreak
\hspace{1em} &  & Yes & 14 & \textbf{0.4499} (0.0127) & \textbf{-1.1827} (0.0205)\\*

\addlinespace[0.5em]
\nopagebreak
\hspace{1em} & Probit link & No & 9 & 0.4578 (0.0098) & -1.1944 (0.0187)\\*
\nopagebreak
\hspace{1em} &  & Yes & 14 & 0.4524 (0.0120) & -1.1844 (0.0225)\\*

\addlinespace[0.5em]
\nopagebreak
\hspace{1em} & Cloglog link & No & 9 & 0.4861 (0.0094) & -1.2653 (0.0164)\\*
\nopagebreak
\hspace{1em} &  & Yes & 14 & 0.4806 (0.0112) & -1.2537 (0.0219)\\*
\cmidrule{1-6}
\addlinespace[0.3em]
\multicolumn{6}{l}{\textbf{\textbf{Discrete-time Markov chain}}}\\*
\hspace{1em} &  & No & -- & 0.4412 (0.0227) & -1.0550 (0.0575)\\*
\cmidrule{1-6}
\addlinespace[0.3em]
\multicolumn{6}{l}{\textbf{\textbf{Markov jump process}}}\\*
\hspace{1em} & $\bm{A}_{\RN{1}}$: Nearest class & No & 4 & 0.4721 (0.0132) & -1.1838 (0.0290)\\*
\nopagebreak
\hspace{1em} &  & Yes (global) & 9 & 0.4648 (0.0150) & -1.1711 (0.0300)\\*
\nopagebreak
\hspace{1em} &  & Yes (state-spec.) & 24 & 0.4352 (0.0147) & -1.0842 (0.0337)\\*

\addlinespace[0.5em]
\nopagebreak
\hspace{1em} & $\bm{A}_{\RN{2}}$: Relaxed Erlang & No & 4 & 0.4528 (0.0150) & -1.0827 (0.0236)\\*
\nopagebreak
\hspace{1em} &  & Yes (global) & 9 & 0.4454 (0.0152) & -1.0721 (0.0195)\\*
\nopagebreak
\hspace{1em} &  & Yes (state-spec.) & 24 & 0.4325 (0.0152) & -1.0649 (0.0244)\\*

\addlinespace[0.5em]
\nopagebreak
\hspace{1em} & $\bm{A}_{\RN{3}}$: Up to 2 classes ahead & No & 7 & 0.4479 (0.0146) & -1.0455 (0.0273)\\*
\nopagebreak
\hspace{1em} &  & Yes (global) & 12 & 0.4406 (0.0149) & -1.0326 (0.0277)\\*
\nopagebreak
\hspace{1em} &  & Yes (state-spec.) & 27 & \textbf{0.4265} (0.0156) & \textbf{-1.0189} (0.0271)\\*

\addlinespace[0.5em]
\nopagebreak
\hspace{1em} & $\bm{A}_{\RN{4}}$: Up to 3 classes ahead & No & 9 & 0.4472 (0.0153) & -1.0389 (0.0280)\\*
\nopagebreak
\hspace{1em} &  & Yes (global) & 14 & 0.4400 (0.0151) & -1.0257 (0.0266)\\*
\nopagebreak
\hspace{1em} &  & Yes (state-spec.) & 29 & 0.4266 (0.0156) & -1.0190 (0.0267)\\*

\addlinespace[0.5em]
\nopagebreak
\hspace{1em} & $\bm{A}_{\RN{5}}$: Free upper-triangular & No & 10 & 0.4472 (0.0153) & -1.0392 (0.0277)\\*
\nopagebreak
\hspace{1em} &  & Yes (global) & 15 & 0.4400 (0.0151) & -1.0260 (0.0262)\\*
\nopagebreak
\hspace{1em} &  & Yes (state-spec.) & 30 & 0.4266 (0.0156) & -1.0191 (0.0267)\\*
\cmidrule{1-6}
\addlinespace[0.3em]
\multicolumn{6}{l}{\textbf{\textbf{Model with modified sojourn time estimates}}}\\*
\hspace{1em} & $\bm{A}_{\RN{1}}$: Nearest class & No & 8 & 0.4584 (0.0149) & -1.1222 (0.0362)\\*
\nopagebreak
\hspace{1em} &  & Yes (global) & 13 & 0.4496 (0.0156) & -1.1047 (0.0350)\\*

\addlinespace[0.5em]
\nopagebreak
\hspace{1em} & $\bm{A}_{\RN{2}}$: Relaxed Erlang & No & 8 & 0.4503 (0.0160) & -1.0797 (0.0261)\\*
\nopagebreak
\hspace{1em} &  & Yes (global) & 13 & 0.4424 (0.0161) & -1.0689 (0.0233)\\*

\addlinespace[0.5em]
\nopagebreak
\hspace{1em} & $\bm{A}_{\RN{3}}$: Up to 2 classes ahead & No & 11 & 0.4446 (0.0145) & -1.0397 (0.0282)\\*
\nopagebreak
\hspace{1em} &  & Yes (global) & 16 & \textbf{0.4371} (0.0152) & -1.0266 (0.0291)\\*

\addlinespace[0.5em]
\nopagebreak
\hspace{1em} & $\bm{A}_{\RN{4}}$: Up to 3 classes ahead & No & 13 & 0.4447 (0.0149) & -1.0374 (0.0276)\\*
\nopagebreak
\hspace{1em} &  & Yes (global) & 18 & 0.4373 (0.0154) & \textbf{-1.0248} (0.0278)\\*

\addlinespace[0.5em]
\nopagebreak
\hspace{1em} & $\bm{A}_{\RN{5}}$: Free upper-triangular & No & 14 & 0.4447 (0.0149) & -1.0376 (0.0274)\\*
\nopagebreak
\hspace{1em} &  & Yes (global) & 19 & 0.4373 (0.0153) & -1.0250 (0.0276)\\*
\cmidrule{1-6}
\addlinespace[0.3em]
\multicolumn{6}{l}{\textbf{\textbf{Mixture Markov jump process ($G=4$)}}}\\*
\hspace{1em} & $\bm{A}_{\RN{1}}$: Nearest class & No & 19 & 0.4400 (0.0206) & -1.0565 (0.0995)\\*
\nopagebreak
\hspace{1em} &  & Yes (global) & 24 & 0.4177 (0.0134) & -0.9716 (0.0313)\\*

\addlinespace[0.5em]
\nopagebreak
\hspace{1em} & $\bm{A}_{\RN{2}}$: Relaxed Erlang & No & 19 & 0.4528 (0.0150) & -1.0827 (0.0236)\\*
\nopagebreak
\hspace{1em} &  & Yes (global) & 24 & 0.4163 (0.0126) & -0.9693 (0.0332)\\*

\addlinespace[0.5em]
\nopagebreak
\hspace{1em} & $\bm{A}_{\RN{3}}$: Up to 2 classes ahead & No & 31 & 0.4209 (0.0141) & -0.9672 (0.0295)\\*
\nopagebreak
\hspace{1em} &  & Yes (global) & 36 & 0.4151 (0.0143) & -0.9555 (0.0285)\\*

\addlinespace[0.5em]
\nopagebreak
\hspace{1em} & $\bm{A}_{\RN{4}}$: Up to 3 classes ahead & No & 39 & 0.4221 (0.0147) & -0.9695 (0.0336)\\*
\nopagebreak
\hspace{1em} &  & Yes (global) & 44 & 0.4145 (0.0137) & -0.9550 (0.0299)\\*

\addlinespace[0.5em]
\nopagebreak
\hspace{1em} & $\bm{A}_{\RN{5}}$: Free upper-triangular & No & 43 & 0.4207 (0.0132) & -0.9658 (0.0290)\\*
\nopagebreak
\hspace{1em} &  & Yes (global) & 48 & \textbf{0.4139} (0.0139) & \textbf{-0.9543} (0.0278)\\*
\cmidrule{1-6}
\addlinespace[0.3em]
\multicolumn{6}{l}{\textbf{\textbf{Ensemble}}}\\*
\hspace{1em} & All & Mixed & -- & 0.4222 (0.0130) & -0.9936 (0.0217)\\*
\end{longtable}



\paragraph{Reference predictors.}
Among the CLMs, the logistic link achieves the best scores, while the complementary log-log link performs worst. Although the logistic assumption underlies the latent variable, this should not be interpreted as evidence that defect propagation itself follows a logistic distribution; only that it provides the most effective link function in this context. 
In comparison, the DTMC performs similar as the logistic link CLM in terms of RPS, but substantially better in terms of the log score.
The uniform distribution performs worst among all investigated predictors (both in terms of RPS and log score), which suggest that the data contain nontrivial information so that more complex models are preferred over the simplest baseline. 


\paragraph{Score differences.}
The MJPs consistently outperform the CLMs in log score, reflecting the advantage in capturing temporal dynamics and state dependence. For RPS, the relative difference is smaller: while MJPs are favored in most cases (and especially for the more complex model specifications), the parameterizations $\bm{A}_{\RN{1}}$ and $\bm{A}_{\RN{2}}$ with either no exogenous information or global regression underperform compared to the logistic CLM. This could suggest that these MJPs are sharper but less well calibrated than the CLM, because the log score is more sensitive to concentration of probability on the realized category, while the RPS rewards good probabilistic performance across the ordinal scale more so than the log score.
However, in almost all cases the two probabilistic scores show the same tendencies in the prediction accuracy across the models, which makes it difficult to comment on deficiencies on sharpness, calibration, and discrimination ability. Also, relative differences in RPS and log score alone are not enough to evaluate these forecast qualities. A detailed diagnostic analysis is therefore provided in \cref{sec:reliability_analysis}.

\paragraph{Generator parameterizations.}
The choice of generator parameterization substantially impacts prediction accuracy. 
The structural differences between the generator parameterizations reflect different assumptions about the physical degradation mechanisms of rail defects.
In summary, $\bm{A}_{\RN{1}}$ enforces gradual progression via neighboring classes, $\bm{A}_{\RN{4}}$ (and similarly $\bm{A}_{\RN{3}}$ and $\bm{A}_{\RN{5}}$) allows for abrupt jumps to severe states (also $\bm{A}_{\RN{2}}$ but more constrained). 
Overall, generator structures that allow direct transitions to non-neighboring classes yield better scores. This suggests that modeling both fatigue-driven crack growth and occasional abrupt worsening due to shocks imposed by rolling stock improves predictive accuracy.
For the standard MJP\footnote{``Standard MJP'' refers to the MJP without modified sojourn time estimates, mixture distributions, either without covariates or else using the global regression setting (\cref{eq:trans_rates_covariates}).}, predictive accuracy increases with generator flexibility, so that $\bm{A}_{\RN{4}}$ and $\bm{A}_{\RN{5}}$ are the best performing parameterizations within that model class. For the more complex model structures (the models with modified sojourn estimates, state-dependent regression coefficients, mixture distributions), the improvements are either smaller or non-existent when moving from $\bm{A}_{\RN{3}}$ to more flexible parameterizations (also supported by the AIC values from Appendix \ref{sec:diagnostics_estimation}).

In practice, selecting a generator parameterization could depend on the characteristics of the physical degradation mechanisms and data characteristics. For instance, if domain knowledge suggests that degradation mostly progresses gradually, neighbor-to-neighbor transitions can be modeled using $\bm{A}_{\RN{1}}$. If sudden damage is expected to happen, then transitions over multiple classes could be allowed in the generator.
If data are not interval-censored (so sojourn times are observable), the generator can be parameterized so it reflects the observed transition behavior. As this is not the case here (and that the exact degradation mechanisms at play in the track are not known), several parameterizations can be benchmarked and selected among using out-of-sample predictive scores, as outlined in this study.

In \cref{sec:results_transient_dist_markov}, we further compare and discuss the transient distributions from selected examples.



\paragraph{Model with modified sojourn time estimates, mixture distributions, and ensemble forecast.}
The model with modified sojourn time estimates (see \cref{sec:warping}) improves prediction accuracy across all parameterizations compared to the standard MJP. This indicates that modifying sojourn time estimates can correct for biases introduced by severity classification schemes, because the effective clock runs faster or slower depending on the occupied class through \cref{eq:warped_time}, which can compensate for defect classifications not being proportional to the physical crack size.

The mixture MJPs yield the best predictive scores when comparing each generator parameterization across the other model specifications. Even though it introduces many additional model parameters to be determined, the mixture MJPs have the lowest AIC values (Appendix \ref{sec:diagnostics_estimation}). The largest improvements are seen in terms of the log score.

The ensemble forecast is among the best-performing predictors, consistent with prior findings that model averaging can improve predictive accuracy \citep{wu2021a}. The structural differences between predictors allow them to capture complementary aspects of the data. By pooling their predictions, ensemble averaging reduces idiosyncratic model errors and can improve calibration. In our experiments, it does not outperform any of the mixture MJPs with covariates.


\paragraph{Impact of introducing exogenous information.}
Across CLM and MJP settings, we see that including exogenous information improves forecast accuracy. This is consistent with \citet{bisgaard2024A}, who found that such covariates explain structural variability between rail lines, supporting the notion that rail networks are heterogeneous and that local conditions at the defect locations must be accounted for, see for example \citet{bai2017a}.
However, this improvement is relatively modest compared with the improvements obtained when moving from one model class to another. For example, the gain from introducing exogenous information into the CLM is smaller than the gain from moving from either CLM type to the DTMC. This suggests that, although defect transition behavior depends on the included covariates, forecast accuracy is improved more by using a model class that better captures the transition behavior in the data. Although there is one caveat to that conclusion: by introducing state-dependent regression coefficients (via \cref{eq:trans_rates_covariates2}), \cref{tab:prediction_scores} shows that predictive accuracy is improved rather substantially across the different MJP generator parameterizations. One reason for this is that \cref{eq:trans_rates_covariates2} incorporates the expected sojourn time estimates $\mu_j^{(s)}$ through the weights $w_j^{(s)}$, which we saw earlier improved the forecast accuracy in the models with modified sojourn time estimates, but even more so for the state-dependent regression setting. What this regression setting also does is to allow covariates to impact transition rates conditional on how progressed defects are, thereby modeling a more complex dependency between transition behavior and covariates. The cost is more parameters, but in this forecasting problem, it improves out-of-sample accuracy despite that.
  

\subsection{Comparing model dynamics} \label{sec:results_transient_dist_markov} 
The investigated models imply substantially different degradation dynamics. In this section, we compare selected predictors to understand how they evolve over time, approximate the observed transition behavior, and the influence of the covariates.

\subsubsection{Transient distributions}

\cref{fig:transient_dists} (page \pageref{fig:transient_dists}) illustrates the transient distributions over a ten-year period for a randomly chosen class $3$ defect
comparing the standard MJP with global regression coefficients (\cref{eq:trans_rates_covariates}) with the state-dependent regression setting (\cref{eq:trans_rates_covariates2}) as well as the mixture MJP (\cref{sec:mixture}). The generator parameterization is held fixed ($\bm{A}_{\RN{3}}$). To the right, the corresponding ten-year transition probability matrices (TPMs) are displayed.


 \begin{figure*}[htb!]
        \centering
        \begin{subfigure}[t]{1 \textwidth}
            \centering
            \includegraphics[width=0.475\textwidth]{trans_dist_gerlang.pdf}
            \includegraphics[width=0.475\textwidth]{tpm_dist_gerlang.pdf}
            \caption[]%
            {{\small Nearest class, that is $\bm{A}_{\RN{1}}$ in \cref{eq:generator_erlang}.}}   
            \label{fig:trans_dist_gerlang}
        \end{subfigure}
        \\ 
        \begin{subfigure}[t]{1 \textwidth}   
            \centering 
            \includegraphics[width=0.475\textwidth]{trans_dist_gerlang_relax.pdf}
             \includegraphics[width=0.475\textwidth]{tpm_dist_gerlang_relax.pdf}
            \caption[]%
            {{\small Relaxed Erlang, that is $\bm{A}_{\RN{2}}$ with mean degradation time following \cref{eq:trans_rate_param}.}}
            \label{fig:trans_dist_gerlang_relax}
        \end{subfigure}
        \\
        \begin{subfigure}[t]{1 \textwidth}  
            \centering 
            \includegraphics[width=0.475\textwidth]{trans_dist_tridiagonal.pdf}
             \includegraphics[width=0.475\textwidth]{tpm_dist_tridiagonal.pdf}
            \caption[]%
            {{\small Up to 3 classes ahead, that is $\bm{A}_{\RN{4}}$ in \cref{eq:generator_tridiagonal}.}}    
            \label{fig:trans_dist_tridiagonal}
        \end{subfigure} \\
        \caption{Left column: Transient distribution over time of the Markov jump process wrt. the parameterization of the generator. The initial distribution of the chain is $\bm{\pi}^{(s)}(0) = \begin{bmatrix}1 & 0 & 0 & 0 & 0 \end{bmatrix}$. Right column: Transition probability matrix, that is, the matrix exponential, $\exp\left(\bm{A}^{(s)}\tau'^{(s)}\right)$, where $\bm{A}^{(s)}$ refers to a chosen parameterization of the generator and $\tau'^{(s)}$ is eight years.}
        \label{fig:transient_dists}
    \end{figure*}

The first row of each TPM corresponds to the transient distribution after ten years. 
The standard MJP places the majority of probability mass in the absorbing class $0$ after ten years (\cref{fig:trans_dist_1}), whereas the MJP with state-dependent regression assigns most probability at class $1$ (\cref{fig:trans_dist_2}), but with substantially less sharpness (i.e. probabilities are more equally spread across classes). The mixture MJP approaches a state distribution which is more static after two years (\cref{fig:trans_dist_3}) than its counterparts.

\paragraph{Turning points in the evolution of class probabilities.} \label{sec:turning_points}
With upper triangular generators, the transient distribution of the standard MJP follows a simple principle: once a class probability starts to decrease, it will not increase again. This results in curves with at most one turning point, cf. \cref{fig:trans_dist_1}. 
This principle does not apply to the MJP with the state-dependent regression of \cref{eq:trans_rates_covariates2}, because as the observation interval increases, the expected proportion of time spent in each class (i.e. the weights $w_j^{(s)}$) changes. For instance, at small observation intervals, most weight will be placed on the starting class, whereas for longer observation intervals more weight will be placed on the more progressed classes.
Since the state-dependent regression depends on these weights by \cref{eq:trans_rates_covariates2} that again depend on the observation interval through \cref{eq:sojourn_time_estimate}, the resulting time-scaling $g(\eta^{(s)}) \tau^{(s)}$ is not necessarily monotonically increasing with $\tau^{(s)}$, because the coefficients $\bm \beta_j$ for every class $j \in \mathcal J$ will dominate according to the weights $w_j^{(s)}$ (so $g(\eta^{(s)})$ can decrease when $\tau^{(s)}$ increases).
Therefore, when the transient distribution is plotted against the observation interval, multiple turning points can occur, which is also evident in \cref{fig:trans_dist_2}. For instance, we see that the class $1$ probability increases rapidly during the first year, then it decreases slowly until year five, where it begins to increase again. Increasing the time horizon would reveal this class $1$ probability to finally approach zero, as the defect reaches the absorbing state.

\paragraph{The implications of different model specifications.}
The models specifications provide different ways to represent the degradation mechanisms. The transient distribution of the standard MJP offers an interpretation of the modeled degradation which is easily conveyable to practitioners: 
It models gradual progression of defects for all of the considered generators,
and it can also model the probability of shock events, by choosing a generator that allows multiple skips.
Comparing the transient distribution of $\bm{A}_{\RN{3}}$ in \cref{fig:trans_dist_1}, with the one that would arise from using $\bm{A}_{\RN{1}}$, one would see substantially more probability mass assigned to class $2B$, because $\bm{A}_{\RN{1}}$ would enforce defects to visit that class before progressing further, even though this class has few observations in the available data (see \cref{fig:hists_u}). Therefore, the standard MJP approach has the ability to adapt to unbalanced data through generator parameterization. 
The model with modified sojourn time estimates reveals the same structural behavior in the transient distribution as the standard MJP, but with higher flexility in how probability mass can be spread according to a modified ``process clock'' instead of calendar time.

The state-dependent regression offers a model extension to the standard MJP that models more complex dynamics between exogenous information and the response in transition behavior. The resulting transient distribution becomes more difficult to communicate, and it may not reflect the ``common sense'' understanding of degradation behavior as seen in \cref{fig:trans_dist_2}, as we described \cref{sec:turning_points}. However, it improves prediction accuracy substantially (cf. \cref{tab:prediction_scores}), compared to the global regression setting, and practitioners could argue for its use based on how well it fits data.

The transient distribution of the mixture MJP constitutes a weighed average of the transient distributions from $G$ MJPs, with transition rates and weights for the latent components jointly determined. With this model complexity, some would argue that it approaches model formulations seen in purely data-driven approaches. It is more difficult to communicate the degradation dynamics, as the latent components may not have a direct physical interpretation, but in the forecasting task it outperforms any other of the predictors (\cref{tab:prediction_scores}).


%The transient distribution is more stable over time,
%(the following section evaluates how well it approximates the transition behavior (independently from observation intervals) in the data).
%complex dependencies with exogeneous information
%how effective it approximates the empirical distribution


%%consider making a new subsection for transition residuals and covariate impact







%Because $\bm{A}_{\RN{1}}$ only allows jumps to neighboring classes, it assigns substantially more probability mass to class $2B$ (see \cref{fig:trans_dist_gerlang}) than $\bm{A}_{\RN{4}}$ (\cref{fig:trans_dist_tridiagonal}), which permits direct jumps to more severe categories without passing through all intermediate states.
%This difference likely reflects the relatively low number of observed class $2B$ defects, causing $\bm{A}_{\RN{4}}$ to assign little probability to that state over time.

%$\bm{A}_{\RN{2}}$ (\cref{fig:trans_dist_gerlang_relax}) assigns less probability mass to class $2B$ than $\bm{A}_{\RN{1}}$, but more than $\bm{A}_{\RN{4}}$. However, it differs notably from both in its treatment of the most severe classes: it assigns substantially more probability mass to class $1$ and correspondingly less to class $0$ after ten years.
%This shift stems from the core assumption of $\bm{A}_{\RN{2}}$ (\cref{eq:trans_rate_param}), where the degradation time between two observed states is modeled as the sum of intermediate sub-degradation processes.
%As a result, probability mass is systematically shifted toward intermediate states, leading to the observed leftward bias in the transition matrix compared to $\bm{A}_{\RN{1}}$ and $\bm{A}_{\RN{4}}$.





%While $\bm{A}_{\RN{1}}$ enforces gradual transitions between neighboring severity classes, it can still capture sudden defect propagation through rapid sequences of small jumps.Conversely, $\bm{A}_{\RN{4}}$ (as well as $\bm{A}_{\RN{3}}$ and $\bm{A}_{\RN{5}}$) allows direct jumps to more severe classes without passing through intermediates, making it well-suited for modeling abrupt failures caused by random shocks, such as those imposed by heavy rolling stock loads or external impacts. $\bm{A}_{\RN{2}}$ lies between these two extremes: it models large transitions as the cumulative effect of smaller degradation steps, combining gradual wear processes with the possibility of accelerated defect growth.

% Note that the first row in the transition probability matrices corresponds to the transient distribution after eight years. As the generalized Erlang only models jumps to the neighboring class, it assigns substantially more probability mass to class $2B$ (see \cref{fig:trans_dist_gerlang}) than the free upper parameterization in \cref{fig:trans_dist_free_upper}, which is not restricted to visit $2B$ before it can move further into more severe categories. Most likely, the relatively low number of $2B$ observations is the reason why the free upper parameterization from \cref{eq:generator_free_upper} assigns low probability to that state over the eight year period. The reparameterized upper in \cref{fig:trans_dist_param_upper} assigns less probability mass to category $2B$ than the generalized Erlang, but more than the free upper. However, the reparameterized upper from \cref{eq:trans_rate_param} differs notably from the two other parameterizations wrt. category $1$ and $0$ where it assigns significantly more probability mass to class $1$ than the generalized Erlang and the free upper. Most of this probability mass comes from class $0$ which has substantially lower probability mass after eight years than in the cases of the generalized Erlang and the free upper. This is caused by the core assumption of the reparameterized upper in \cref{eq:trans_rate_param}, where the degradation time between two observed states is the sum of all intermediate sub-degradation processes. In general, the assumption will force more probability mass to the left in the transition probability matrix, which indeed is the case in \cref{fig:trans_dist_param_upper} compared to \cref{fig:trans_dist_gerlang} or \cref{fig:trans_dist_free_upper}.
% Although, the free upper parameterization has the highest prediction accuracy, the reparameterized upper has its merits. The parameter spaces of the free upper increases quadratically with the number of categories in the state space, while it increases linearly for the reparameterized upper. For the fairly small state space considered in this present study ($m = 5$), model estimation appears robust. However, for larger state spaces it might be more problematic to robustly determine all parameters using the free upper parameterization. In that case, the reparameterized upper could offer an alternative with as few parameters as the generalized Erlang but with more flexibility for jumps in the Markov chain.

%\FloatBarrier

%\subsection{Exogenous information} \label{sec:results_exogenous_info}
%The inclusion of the exogenous information containing the track characteristics consistently improves prediction errors. This holds true for all investigated model settings of the cumulative link model and the Markov jump process. 
%This supports the notion of heterogeneity of the rail network, since taking the specific rail characteristics at the defect location into account leads to better predictions, see \cref{tab:forecast_error_tab}.
%\citet{bisgaard2024A} demonstrate that the included covariates are statistically significant on a 5\% level, computing the standard errors of the parameters via a second order approximation of the score function in the estimated optimum. Due to the $k$-fold cross validation procedure in this study, we obtain $k$ different estimates for each variable for each model setting. In \cref{tab:covariates_estimates}, we show the range of the coefficients.
%\begin{table}[htb!]
%    \centering
%    \caption{Model estimates of rail track characteristics from k-fold estimation procedure}
%    \centerline{\begin{tabular}[t]{>{}llccccc}
\toprule
 & Model & Tonnage & Line speed & Rail profile & Steel hardness & Curvature\\
\midrule
\addlinespace[0.3em]
\multicolumn{7}{l}{\textbf{\textbf{Cumulative link}}}\\
\hspace{1em} & Logistic link & {}[0.49, 0.66] & {}[0.03, 0.15] & {}[-1.15, -0.95] & {}[0.80, 0.91] & {}[0.48, 0.68]\\

\hspace{1em} & Probit link & {}[0.33, 0.42] & {}[0.01, 0.08] & {}[-0.72, -0.59] & {}[0.44, 0.49] & {}[0.28, 0.39]\\

\hspace{1em} & Cloglog link & {}[0.40, 0.50] & {}[-0.10, 0.01] & {}[-0.86, -0.68] & {}[0.42, 0.48] & {}[0.27, 0.40]\\
\cmidrule{1-7}
\addlinespace[0.3em]
\multicolumn{7}{l}{\textbf{\textbf{Markov jump process}}}\\
\hspace{1em} & Generalized Erlang, $\bm{A}'$ & {}[0.38, 0.49] & {}[0.30, 0.42] & {}[-0.86, -0.65] & {}[0.69, 0.78] & {}[0.45, 0.59]\\

\hspace{1em} & Reparameterized upper, $\bm{A}''$ & {}[0.57, 0.72] & {}[0.20, 0.31] & {}[-0.78, -0.58] & {}[0.69, 0.78] & {}[0.52, 0.65]\\

\hspace{1em} & Free upper, $\bm{A}'''$ & {}[0.57, 0.69] & {}[0.15, 0.25] & {}[-0.61, -0.41] & {}[0.65, 0.74] & {}[0.58, 0.71]\\
\bottomrule
\end{tabular}}
%    \label{tab:covariates_estimates}
%\end{table}
%For all folds, the coefficients on the covariates have constant sign and are of similar magnitude, with the only exception of the complementary loglog link for the cumulative link model. However as shown in \cref{tab:forecast_error_tab}, the complementary loglog link has the worst prediction performance of all investigated parametric models, suggesting that the assumption of a min-Gumbell distributed latent variable is not suitable. With that aside, the fact that different settings (link functions and generator parameterizations) of two distinct statistical models lead to similar estimated coefficients for the covariates on all studied subsets of the data strongly suggests that the estimations are robust and that the statistical interpretation of the influence of the covariates is sound. We conclude that the models associate increases in tonnage, line speed, steel hardness and curvature with faster propagation of defect severity. Increasing rail profile is associated with slower propagation of defect severity. As argued by \citet{bisgaard2024A}, these associations are in accordance with railway domain knowledge \citep{UIC_2002_defecthandbook,RASMUSSEN20178}


\subsubsection{Transition residuals} \label{sec:transition_patterns}

To compare the model-implied transitions with the empirical ones, define the elements in the transition count matrix
$
  N_{i,j}^{\text{obs.}} = \sum_{s \in \mathcal S} \mathds{1}_{y_1^{(s)} =\; i, y_2^{(s)} =\; j}
$ for $i,j \in \mathcal J$, which express the number of observed transitions from state $i$ to $j$. The expected transition counts implied by a model predicting a transient distribution $\widehat{\bm \pi}^{(s)}$ at the time of the second observation is given by
$
\widehat{  N}_{i,j} =  \sum_{s \in \mathcal S} \mathds{1}_{y_1^{(s)} = \; i} \widehat{ \pi}_j^{(s)},
$ for $i,j \in \mathcal J$.
Then, the transition residuals for a model can be calculated as
$$
  R_{i,j} = \left(  N_{i,j}^{\text{obs.}} - \widehat{  N}_{i,j} \right) / \sum_{j}   \widehat{  N}_{i,j},
$$ for $i,j \in \mathcal J$.
The entries, $  R_{i,j}$, express the error between the observed number of transitions from state $i$ to $j$, and the expected counts implied by the model predictions, in proportions. 
By construction, $\sum_{j}   R_{i,j} = 0$, so if $  R_{i,j} > 0$, the model underpredicts the amount of $i$ to $j$ transitions, and vice versa for $  R_{i,j} < 0$.
Thus, inspecting $\bm R$ (comprised by the elements $R_{i,j}$) for a model can serve as check of how well it approximates the empirical distribution of the transition patterns between starting and ending state. The norm of $\bm R$ can be used as simple diagnostic for this check, e.g. taking the L1-norm of the transition residuals
$
|| \bm R || = \sum_{i,j} |  R_{i,j}|
$ for $i,j \in \mathcal J$.

For each predictor, we calculate $\bm R$ based on the held-out set in the CV described in \cref{sec:prediction_error}, then aggregating across the folds. In \cref{tab:transition_patterns}, three of these $\bm R$'s are displayed for the MJP under the $\bm{A}_{\RN{3}}$ parameterization with the different covariate setting in \cref{sec:regression} (\cref{tab:expected_transition_bidiagonal_cov1_exp_exp_all_no_warp}, respectively \cref{tab:expected_transition_bidiagonal_cov2_exp_exp_all_no_warp}), and a mixture model using four generators (\cref{tab:expected_transition_bidiagonal_cov1_exp_exp_all_no_warp_mixture4}). As a reference check of the empirical distribution of the $i$ to $j$ transitions, the empirical row-normalized transition matrix, $ N_{i,j}^{\text{obs.}} / \sum_{j}  N_{i,j}^{\text{obs.}}$, is included (\cref{tab:transition_matrix}).

\begin{table}[h]
    \begin{subtable}[h]{0.45\textwidth}
        \centering
            \input{transition_matrix}
          \caption{Transition matrix, $ N_{i,j}^{\text{obs.}} / \sum_{j}  N_{i,j}^{\text{obs.}}$.}
       \label{tab:transition_matrix}
    \end{subtable}
    \hfill
    \begin{subtable}[h]{0.45\textwidth}
        \centering
       \input{expected_transition_bidiagonal_cov1_exp_exp_all_no_warp}
        \caption{$\bm R$ given $\bm{A}_{\RN{3}}$ with global regression (\cref{eq:trans_rates_covariates}).}
        \label{tab:expected_transition_bidiagonal_cov1_exp_exp_all_no_warp}
     \end{subtable}
      \begin{subtable}[h]{0.45\textwidth} 
        \centering
         \input{expected_transition_bidiagonal_cov2_exp_exp_all_no_warp}
        \caption{$\bm R$ given $\bm{A}_{\RN{3}}$ with state-dependent regression (\cref{eq:trans_rates_covariates2}).}
        \label{tab:expected_transition_bidiagonal_cov2_exp_exp_all_no_warp}
     \end{subtable}
     \hfill 
     \begin{subtable}[h]{0.45\textwidth} 
        \centering 
 \input{expected_transition_bidiagonal_cov1_exp_exp_all_no_warp_mixture4}
        \caption{$\bm R$ given mixtures of $\bm{A}_{\RN{3}}$ (\cref{sec:mixture} with $G = 4$) and global regression (\cref{eq:trans_rates_covariates}).}
        \label{tab:expected_transition_bidiagonal_cov1_exp_exp_all_no_warp_mixture4}
     \end{subtable} 
     \caption{Transition residuals comparing model-implied and empirical transitions.}
     \label{tab:transition_patterns}
\end{table}


As seen from the proportions in \cref{tab:transition_matrix}, there is a substantial amount of transitions appearing over more than one class ahead. Overall, the $\bm{A}_{\RN{3}}$ with state-dependent covariates (\cref{tab:expected_transition_bidiagonal_cov2_exp_exp_all_no_warp}) improves the residuals compared to the setting with global regression coefficients (\cref{tab:expected_transition_bidiagonal_cov1_exp_exp_all_no_warp}), as indicated by a lower L1-norm (0.385 vs 0.327). Both settings have their largest residuals in the $1 \rightarrow 0$ transition and overpredict the probability of reaching the absorbing state (class $0$). Clearly, the mixture model matches the empirical transition frequencies more closely (\cref{tab:expected_transition_bidiagonal_cov1_exp_exp_all_no_warp_mixture4}), with an L1-norm of approx. 0.020, which is expected given its flexibility in adapting to data. In contrast, the single-generator MJP settings are more constrained by the assumed degradation mechanism, while the observed transition data may be influenced by practices and policies that prevent many defects from progressing to rail breaks once they reach a safety-critical level (e.g., class $1$).

Appendix \ref{sec:transition_L1} provides a full summary of the L1-norms of the $\bm R$'s from all outlined predictors.




\subsubsection{The impact of the covariates: An example} \label{sec:risk_curves}
In practice, infrastructure managers could be interested in quantifying the impact of the covariates on the risk of obtaining a critical defect. We here provide an example of how this can be done within the existing model framework.
Define the critical class set $\mathcal E_{\text{critical}} = \{ 1,0\}$ with corresponding index set $\mathcal J_{\text{critical}}$. Further, we define the associated critical curve
$$
r^{(s)}(t) = \mathds P\left(X^{(s)}(t) \in \mathcal J_{\text{critical}} \mid X^{(s)}(0) = y_1^{(s)}, \bm z^{(s)}\right) = \sum_{j\in \mathcal J_{\text{critical}}} \widehat{\pi}_j^{(s)}(t; \bm z^{(s)}), 
$$
where $\widehat{\bm \pi}^{(s)}(t; \bm z^{(s)})$ is the predicted distribution for defect $s \in \mathcal S$ at time $t$ with covariates $\bm z^{(s)}$.

To exemplify how covariates impact the risk curve $r^{(s)}(t)$, we consider an MJP with a regression setting. For this example, we use the $\bm{A}_{\RN{3}}$ generator with state-dependent regression coefficients (\cref{eq:trans_rates_covariates2}) and set $y_1^{(s)}$ equal to class $3$. Three scenarios are then computed for a single covariate $\bm z_k \in \bm z$, that is the empirical 10\%, 50\%, and 90\% quantiles across the defects in the dataset, while the empirical means are computed for the remaining covariates. 
For each scenario, we then compute $r^{(s)}(t)$ over a five-year horizon for the variables tonnage, curvature, and rail weight. These risk curves are presented in \cref{fig:critical_curves}, so that covariate-dependency is reflected through the quantiles.

\begin{figure}[htb!]
  \centering
\begin{subfigure}{.33\textwidth}
  \centering
  \includegraphics[width=\linewidth]{crit_curves_1}
  \caption{Tonnage}
  \label{fig:crit_curves_1}
\end{subfigure}%
\begin{subfigure}{.33\textwidth}
  \centering
  \includegraphics[width=\linewidth]{crit_curves_5}
  \caption{Curvature}
  \label{fig:crit_curves_5}
\end{subfigure}
\begin{subfigure}{.33\textwidth}
  \centering
  \includegraphics[width=\linewidth]{crit_curves_3}
  \caption{Rail weight}
  \label{fig:crit_curves_3}
\end{subfigure}
\caption{Risk curves $r^{(s)}(t)$ under varying impacts of covariates using an MJP with $\bm{A}_{\RN{3}}$ generator and state-dependent regression coefficients (\cref{eq:trans_rates_covariates2}).}
\label{fig:critical_curves}
\end{figure}

From \cref{fig:crit_curves_1} and \cref{fig:crit_curves_5}, we see that the model associates high tonnage and curvature values with higher risk of reaching $\mathcal J_{\text{critical}}$, in line with existing domain knowledge that increasing intensity of traffic increases defect progression. \cref{fig:crit_curves_3} indicates that bigger rail profiles lower the risk of a defect progression into a critical class.
In summary,  the resulting curves provide a simple, decision-relevant visualization of how changes in covariates shift the predicted probability of reaching a critical defect class (given a model specification) over time. 


\subsection{Calibration, discrimination, and sharpness} \label{sec:reliability_analysis}
The proper scoring rules in \cref{sec:prediction_error} showed that the flexible MJP parameterizations outperform the reference predictors (most substantially the mixture MJPs), and that modifying sojourn time estimates, exogenous information, and averaging over all predictors in an ensemble improve accuracy.
However, such scores reduce forecast quality to single summary values, and therefore cannot reveal the specific sources of deficiencies \citep{hamill1997,Dimitriadis_2024}. So far, we supplemented this evaluation by inspecting transition residuals (\cref{sec:transition_patterns}), but to provide a more detailed analysis, multicategorical reliability diagrams are generated to assess forecast calibration, and checkerboard plots to evaluate discrimination ability.\footnote{Our GitHub repository (\url{https://github.com/dobbeltgaard/MarkovJump.git}) provides \texttt{R} code to generate these diagnostic plots.} Also, the sharpness of the predictors will be assessed using an entropy measure. 

\subsubsection{Multicategory reliability diagrams}
We follow the methodology proposed by \citet{hamill1997}, extending the framework of reliability diagrams initially introduced for binary forecasts by \citet{wilks1995a}.
The diagnostic plots are generated by the following procedure.
\begin{enumerate}    
    \item Define a set of quantiles $q_1, q_2, \dots, q_K$ (e.g., $0.05, 0.15, \dots, 0.95$), where $K$ is the number of quantiles.
    \item For each forecast, compute the cumulative probabilities across the $m$ categories. For each quantile $q_k$, determine the forecasted category as the smallest category where the cumulative probability exceeds or equals $q_k$. This converts the probabilistic forecasts into categories at each quantile.
    \item For each forecast-observation pair, the multicategory reliability is defined as:
    \begin{itemize}
        \item $1$ if the observed category is less than the forecasted category.
        \item $0$ if the observed category is greater than the forecasted category.
        \item $\frac{q - q_{\min}}{q_{\max} - q_{\min}}$ if the observed category equals the forecasted category, where $q_{\min}$ and $q_{\max}$ are the quantiles where the observation matches the forecast.
    \end{itemize}
    \item For each quantile $q_k$, compute the average reliability across all forecast-observation pairs.
\end{enumerate}
As a supplementary summary metric, we report the \textit{average category error}, defined as the mean absolute difference between the forecasted category at each quantile and the observed category, averaged first across forecasts and then across quantiles.

Error bars are calculated via a bootstrap procedure, which repeats the above steps across multiple resampled datasets. A checkerboard plot of the difference between the forecasted and observed categories at each quantile is used to visualize the discrimination ability of the forecasts. \citet{hamill1997} offers an exhaustive description of the above methodology.

\subsubsection{Entropy measure}
To assess forecast sharpness, we define for a predicted distribution $\widehat{\bm\pi}^{(s)}$ a normalized entropy measure by
\begin{align}
H(\widehat{\bm\pi}^{(s)}) = - \frac{1}{\log m}\sum_{j=1}^{m}\widehat{\pi}^{(s)}_j \log\!\big(\widehat{\pi}^{(s)}_j\big), \quad s \in \mathcal S.\label{eq:entropy}
\end{align}
This quantifies the average level of uncertainty that a prediction assigns over the categories with $H \in [0,1]$, so that the sharper a forecast is, the lower $H$ becomes.

\subsubsection{Summary of calibration, discrimination, and sharpness}
Based on the out-of-fold predictions from the CV procedure in \cref{sec:prediction_error}, the multicategory reliability diagrams with checkerboard plots are generated, so each defect contributes once. Further, we report mean predictive entropy as a measure of sharpness, computed for each held-out forecast and averaged (a full list is provided in Appendix \ref{sec:sharpness}).

To illustrate, we focus on four predictors: 
1) the uniform predictor; 
2) the logistic CLM with exogenous information; 
3) the MJP with parameterization $\bm{A}_{\RN{3}}$ including exogenous information;
4) and the MJP mixture using the $\bm{A}_{\RN{3}}$ parameterization including exogenous information. 
Bootstrapping was performed with 200 resamples, and 10\% and 90\% confidence intervals are reported. We use $K = 10$ quantiles ($0.05, 0.15, \dots, 0.95$). Results are shown in \cref{fig:reliability_diagrams}.
Note that wider uncertainty bands at the highest quantiles can reflect lower effective sample sizes, particularly for rare class $0$ transitions. We found 200 resamples sufficient for stable curves in repeated runs.



 \begin{figure*}[htb!]
        \centering
        \begin{subfigure}[t]{0.475\textwidth}
            \centering
            \includegraphics[width=\textwidth]{reliability_diag_uniform.pdf}
            \caption[]%
            {{\small Uniform distribution predictor from \cref{eq:uniform_predictor}}}    
            \label{fig:reliability_diag_uniform}
        \end{subfigure}
        \hfill
        \begin{subfigure}[t]{0.475\textwidth}  
            \centering 
            \includegraphics[width=\textwidth]{reliability_diag_olr_cov.pdf}
            \caption[]%
            {{\small Cumulative logistic link with exogenous information}}    
            \label{fig:reliability_diag_olr_cov}
        \end{subfigure}
        \hfill
        \begin{subfigure}[t]{0.475\textwidth}   
            \centering 
            \includegraphics[width=\textwidth]{reliability_diag_bidiagonal_TRUE_exp_exp_all_warp.pdf}
            \caption[]%
            {{\small Up to 2 classes ahead with exogenous information and modified sojourn time estimates, that is $\bm{A}_{\RN{3}}$ in \cref{eq:generator_bidiagonal} using \cref{eq:trans_rates_covariates} and \cref{eq:warped_time}.}}    
            \label{fig:reliability_diag_bidiagonal_TRUE_exp_exp_all_warp}
        \end{subfigure}
        \hfill
        \begin{subfigure}[t]{0.475\textwidth}   
            \centering 
            \includegraphics[width=\textwidth]{reliability_diag_ensemble_all.pdf}
            \caption[]%
            {{\small Ensemble}}    
            \label{fig:reliability_diag_ensemble_all}
        \end{subfigure}
        \caption{Multicategorical reliability diagrams showing the fraction of observations below the forecast quantile. We use 10 quantiles, that is ($0.05, 0.15, \dots, 0.95$). The dotted line represents a perfectly calibrated forecast. Error bars are calculated using bootstrap testing. Inset checkerboard plots show the average category error wrt. the quantiles. Highly populated bins are darker colored. Perfect forecasts have 0 category error.}
        \label{fig:reliability_diagrams}
    \end{figure*}


The uniform distribution predictor (\cref{fig:reliability_diag_uniform}) substantially overestimates the fraction of observed categories across all quantiles, indicating high miscalibration.
Its average category error is also the highest among the four predictors, consistent with its poor performance on aggregate scores in \cref{tab:prediction_scores}. Its average normalized entropy is $0.8093$, which is expected since a uniform predictor is preferenceless, meaning it has minimal sharpness.\footnote{Remark the only reason it does not attain the maximum value of $1$ for the average entropy is that \cref{eq:uniform_predictor} only spreads out probability mass to the states that are reachable given the initial observation $y_1^{(s)}$ (if the support size is $r<m$, the maximum normalized entropy becomes $\log(r)/\log(m) < 1$).}

The cumulative link model with logistic link (\cref{fig:reliability_diag_olr_cov}) shows improvement in calibration relative to the uniform distribution predictor.
Its discrimination ability remains modest, as indicated by large average category errors and checkerboard plots with heavily populated bins far from the zero-error line. Average entropy is large ($0.7416$).

The MJP (\cref{fig:reliability_diag_bidiagonal_cov1_exp_exp_all_no_warp}) achieves the lowest mean category error among the single predictors, and its checkerboard plot shows concentration of errors near the zero-category line, indicating good discrimination ability.
Its reliability curve indicates that it tends to underpredict severity, most noticeably at higher quantiles.
The mixture MJP (\cref{fig:reliability_diag_bidiagonal_cov1_exp_exp_all_no_warp_mixture4}) has larger mean category error, but better calibrated forecasts at the higher quantiles (cf. \cref{sec:transition_patterns}, where it better matches empirical transition frequencies). However, it has weakened discrimination ability at the higher quantiles (see the darker bins away from zero at the two last bins in the checkerboard), indicating a trade-off between calibration and discrimination ability.
The entropy of the single generator MJP is $0.5249$ compared to $0.5857$ for the mixture MJP, which suggests that the mixture MJP is more conservative in how it distributes the probability mass (also reflected in slower evolution of the rail break risk in \cref{sec:results_transient_dist_markov}). As the mixture MJP was the predictor with the highest predictive accuracy in terms of RPS and log score, it suggests that a loss of sharpness can be beneficial if it improves calibration enough.



\subsection{Proposition for use case} \label{sec:proposition_for_use}
From the perspective of an infrastructure manager, probabilistic forecasts in isolation have limited operational value. 
For better network operation, forecasts must be embedded in decision frameworks that allocate resources under uncertainty.\footnote{Discussions with BaneNOR and Banedanmark (the infrastructure managers responsible for most of the Norwegian and Danish networks) highlighted that connecting predictions to actionable maintenance strategies is essential for realizing the value of predictive models.}
%One possible implementation pathway could be threshold-based decisions under a fixed budget: for each defect, generate forecasts over a six-month horizon and allocate resources to the $n$ defects with the highest probability of reaching a critical state, with $n$ bounded by the budget divided by the average per-defect action cost. Such a heuristic is unlikely to lead to near-optimal network operation, because it ignores the value of information from inspection and the timing of decisions over multiple horizons.
In the following, we outline and solve a stochastic program that connects probabilistic forecasts to sequential inspection and maintenance decisions over multiple stages. This example is intended as a proof-of-concept rather than a fully operational maintenance planning tool (cf. \cref{sec:conclusion}).


\subsubsection{A stochastic program}
Here, we formulate an optimization model that, based on the predictions, minimizes expected cost over a time horizon of $T_n$ related to active defects in a set $\mathcal{S}$ across a railway network with a rail line mapping $\ell:\mathcal{S}\to\{\text{rail lines}\}$.


\paragraph{Decision variables and available information.}
We will consider maintenance and inspection planning on a time grid $\widehat T=(T_1,\dots,T_n)$.
At these time points (stages), we define decision variables 
$y^{(s)}(t)\in\{0,1\}$ for inspection and $w^{(s)}(t)\in\{0,1\}$ for maintenance at time $t \in \widehat T$ for defect $s\in\mathcal S$.
%$$
%y^{(s)}(t)\in\{0,1\}\quad\text{(inspection at \(T_t\))}, \qquad %\text{and }
%z^{(s)}(t)\in\{0,1\}\quad\text{(maintenance at \(T_t\))},
%$$
%for $t = 1, \dots, n$.
The information available to the infrastructure manager is represented by the $\sigma$-algebra generated by the outcomes of decisions made so far, that is, a filtration
$
\mathcal F(t) = \sigma \left( O^{(s)}(\tau), y^{(s)}(\tau), w^{(s)}(\tau) : \tau \leq t, s\in\mathcal S \right), 
$
for $t = 1, \dots, n$ with the observation process 
$$
O^{(s)}(t) = \begin{cases}
X^{(s)}(t), \quad &y^{(s)}(t) = 1\\ 
\varnothing, \quad &y^{(s)}(t) = 0.
\end{cases}
$$
Note that at each stage $t$, the predicted distributions $\widehat{\bm \pi}^{(s)}(t \mid \mathcal F(t))$ are computed given the information available in $\mathcal F(t)$, and so they are $\mathcal F(t)$-measurable.

We let the indicator \(a^{(s)}(t)\in\{0,1\}\) express that defect \(s\) is still active, which means not maintained yet. We define \(\chi^{(s)}(t)\in\{0,1\}\) as an indicator that defect \(s\) is still uninspected. Subsequently, $a^{(s)}(t)$ and $\chi^{(s)}(t)$ are auxiliary variables that will help us write up the optimization program compactly.


\paragraph{Costs and critical class risks.}
Let the inspection cost at line $\ell$ be denoted by $c_{\text{ins}}^{\ell(s)}$, and $c_{\text{fail}}^{\ell(s)}$ the estimated costs of a rail break for $s \in \mathcal S$.
The expected maintenance cost is defined as 
$
\bar c^{(s)}(t) = \mathds E \left[ c_{\text{maint}}^{\ell(s)} \left(X^{(s)}(t)\right) \mid \mathcal F(t) \right]
$ where $c^{\ell(s)}_{\text{maint}}(j)$ is the maintenance cost for severity class $j$. The planned budget size is denoted $B(T_k)$ for period $T_k$.

Further, we define a critical-state probability, i.e., 
$
q^{(s)}(t) = \mathds P\left(X^{(s)}(t) \in \mathcal J_{\text{critical}} \mid \mathcal F(t) \right)
$ 
where $\mathcal J_{\text{critical}}$ for example can be defined as in \cref{sec:risk_curves}.
Before inspection, $\bar c^{(s)}(t)$ and $q^{(s)}(t)$ are forecast-based; after inspection, they are conditional on the observed class.

\paragraph{The objective function and constraints.}
The multistage stochastic program can then be written as a minimization of the expected costs, that is,
\begin{align}
\min_{y,w,a,\chi}\quad
\mathbb E\!\left[
\sum_{k=1}^n\sum_{s\in\mathcal S}
\Bigl(
c_{\mathrm{ins}}^{\ell(s)} y^{(s)}(T_k)
+ \bar c^{(s)}(T_k) w^{(s)}(T_k)
+ c_{\mathrm{fail}}^{\ell(s)}\,q^{(s)}(T_k)\bigl(a^{(s)}(T_k)-w^{(s)}(T_k)\bigr)
\Bigr)
\right]
\label{eq:general_textbook_obj}
\end{align}
subject to
\begin{align}
\sum_{s\in\mathcal S}
\Bigl(
c^{\ell(s)}_{\mathrm{ins}} y^{(s)} (T_k)
+ \bar c^{(s)}(T_k) w^{(s)}(T_k)
\Bigr)
&\le B(T_k),\quad k = 1,\dots,n,
\label{eq:general_textbook_budget}\\
y^{(s)}(T_k)+w^{(s)}(T_k)
&\le a^{(s)}(T_k),\qquad s\in\mathcal S,\  k = 1,\dots,n,
\label{eq:general_textbook_exclusive}\\
y^{(s)}(T_k)
&\le \chi^{(s)}(T_k),\qquad s\in\mathcal S,\  k = 1,\dots,n,
\label{eq:general_textbook_inspection}\\
a^{(s)}(T_1)&=1,\qquad \chi^{(s)}(T_1)=1, \qquad s\in\mathcal S,
\label{eq:general_textbook_initial}\\
a^{(s)}(T_{k+1})&=a^{(s)}(T_k)-w^{(s)}(T_k),\qquad s\in\mathcal S,\ k=1,\dots,n-1,
\label{eq:general_textbook_activity}\\
\chi^{(s)}(T_{k+1})&=\chi^{(s)}(T_k)-y^{(s)}(T_k),\qquad s\in\mathcal S,\ k=1,\dots,n-1.
\label{eq:general_textbook_uninspected}
\end{align}
%All decision variables and auxiliary state variables are required to be \(\mathcal F_t\)-measurable at stage \(t\). 
The objective in \cref{eq:general_textbook_obj} combines inspection cost, expected maintenance cost, and expected failure cost. 
%Note that $q^{(s)}(t)$ lumps the forecast distribution into a stage-wise critical risk that enters as expected failure exposure for active defects. 
\cref{eq:general_textbook_budget}--\cref{eq:general_textbook_inspection} ensure budget feasibility\footnote{The budget constraint is imposed in expectation, so it should here be interpreted as a planning budget constraint rather than a hard realized budget constraint.} and at most one action per active defect and stage.
\cref{eq:general_textbook_initial}--\cref{eq:general_textbook_uninspected} ensure that the auxiliary variable are updated, so that maintenance removes a defect from the future decision stage and inspections update the information available for later decisions.%\footnote{If an event has probability 1, then we say that this event occurs ``almost surely'' (a.s.) \citep[Chapter 3]{thygesen2023a}.}
Note that this model formulation can be extended to include further operational constraints, such as crew availability, possession windows, travel logistics, etc.


\subsubsection{Solving a stochastic program: An example}
To solve a stochastic program like \cref{eq:general_textbook_obj}--\cref{eq:general_textbook_uninspected}, uncertainty is typically represented by generating scenarios structured in a decision tree.
Given $m$ severity classes, $n$ decision stages, and $|\mathcal S|$ defects, the number of latent degradation paths is in the order of $m^{n|\mathcal S|}$.
Because the information available to the decision-maker depends on what is revealed by inspections, the optimization problem is subject to \textit{endogeneity}, 
%the uncertainty is \textit{endogenous}, 
meaning that the filtration $\mathcal F(t)$ depends on the inspection decisions $y^{(s)}(\tau)$ for $\tau\le t$, which increases the number of information branches.
This means that the optimization problem very quickly becomes intractable for increasing numbers of defects and decision stages.
Therefore, we have to approximate the full decision tree by a smaller one.


\paragraph{Numerical solution.}
First, for $10$ selected defects we generate Monte Carlo scenarios on a three-stage grid $(T_1,T_2,T_3)$ with known severity classes at $T_0$.\footnote{First-stage decisions will therefore be made at $T_1$.}
This is done over a $1.5$ year horizon using transition probabilities from a fitted MJP with $\bm{A}_{\RN{3}}$-generator and regression (\cref{eq:trans_rates_covariates}). 
Second, we reduce these sampled scenarios by using stage-wise clustering of the joint states at $T_2$ and $T_3$.
Endogenous information is handled by letting an inspection reveal the realized class at the decision stage. In later stages, the quantities $q^{(s)}(t)$ and $\bar c^{(s)}(t)$ are then computed conditional on the observed class; if no inspection is made, they are computed from the forecast distributions. More details are provided in Appendix \ref{sec:stochastic_optimization_details}.
\cref{tab:decisions} displays the obtained first-stage decisions.
\begin{table}[htb!]
\centering
\caption{First-stage decisions in the obtained policy.} \label{tab:decisions}
\begin{tabular}{clcl}
\toprule
Defect & Class at $T_0$ & \(q_{\mathrm{break}}(T_1)\) & Stage-1 decision \\
\midrule
1 & Class $3$ & 0.0346 & Inspect \\
2 & Class $2A$ & 0.1564 & Maintain \\
3 & Class $3$ & 0.0285 & Inspect \\
4 & Class $2A$ & 0.0677 & Wait \\
5 &Class $3$ & 0.0165 & Inspect \\
6 & Class $3$ & 0.0286 & Wait \\
7 & Class $2A$ & 0.0912 & Wait \\
8 & Class $2B$ & 0.0747 & Maintain \\
9 & Class $2A$ & 0.0993 & Maintain \\
10 & Class $2A$ & 0.1182 & Maintain \\
\bottomrule
\end{tabular}
\end{table}
As seen, the defects with lower severity are either assigned an inspection or no action, while the more progressed defects are assigned maintenance.


Taking defects 5 and 6 as a small case study, we see that defect 5 is inspected while defect 6 is not, even though both start in class $3$. The normalized entropy (\cref{eq:entropy}) of the one-step forecast at $T_1$ is $0.48$ for defect 5 and $0.66$ for defect 6, so defect 6 has the less sharp forecast at the first stage. %However, the policy is not driven by forecast uncertainty alone: inspection is only valuable if it is likely to change downstream actions.
In our solved policy (not fully shown in \cref{tab:decisions}), defect 6 is scheduled for maintenance at $T_2$ in a large fraction of the scenarios (about $71\%$), whereas defect 5 is maintained at $T_2$ much less often (about $34\%$). In other words, defect 6 is ``almost maintained anyway'', so paying for an inspection rarely changes the decision path. This illustrates how the program trades off inspection cost against the expected value of information.


%> entropy(t(c(1,0,0,0,0)) %*% TPMS[[3]]$TPM$dt_0.5)
%[1] 0.6570918
%> entropy(t(c(1,0,0,0,0)) %*% TPMS[[5]]$TPM$dt_0.5)
%[1] 0.4813692
%> entropy(t(c(1,0,0,0,0)) %*% TPMS[[6]]$TPM$dt_0.5)
%[1] 0.6575578

%questions for marcel: 
%root class or initial distribution?


\section{Discussion and conclusions} \label{sec:conclusion}
We have proposed modeling partially observed rail defect severity trajectories using a Markov jump process (MJP) estimated via optimum score estimation to align estimation with the chosen forecast evaluation metric. The framework accommodates covariate-dependence (tonnage, line speed, curvature, grade, rail profile) through a global regression and one with state-dependent regression coefficients. The study also considers multiple generator parameterizations, a modification of sojourn time estimates, and an extension to a mixture model.
Some of the model specifications are paired, but not all (e.g., state-dependence regression is not combined with mixture MJPs, and mixture MJPs incorporate the same generator parameterization for all latent components). This remains a limitation to this study that investigates model specifications somewhat independently from each other. On the other hand, this approach enables direct comparisons of what introduced model complexity that yields the largest performance improvement.

Through systematic forecast evaluation with proper scoring rules, we benchmarked MJPs against reference predictors, including a uniform predictor, a discrete-time Markov chain, and cumulative link models with logistic, probit, and complementary log-log links. Results from $k$-fold cross-validation show that MJPs consistently outperform all reference predictors, achieving higher accuracy across both log score and ranked probability score; however, estimation of prediction error through $k$-fold cross-validation can be biased \citep{fushiki2011a}.

The choice of generator parameterization affects predictive performance: more flexible structures yield higher accuracy in our experiments. Modifying sojourn time estimates improves performance across all parameterizations. 
The mixture MJPs perform best across all predictors in terms of predictive accuracy. Forecast diagnostics reveal that the mixture MJP approximates the transition behavior in the data well, while they tend to be less sharp than the single-generator MJPs.

Including exogenous information results in increased predictive accuracy across all regression-based predictors (which strongly suggests that accounting for the heterogeneity of local track conditions is beneficial). We see further improvements from introducing state-dependency in the regression, which may indicate complex dependencies at play between transition behavior and covariates. This study's investigation of which covariates are the most important is somewhat limited, as it builds upon the work of \citet{bisgaard2024A}, and most importantly, due to data availability. We encourage future work to provide insight into the statistical associations between crack dynamics and other track related measurements.

Finally, we showed how predictive distributions can be embedded in a stochastic optimization program to allocate inspection and maintenance actions under budget constraints, minimizing expected costs. 
Since the decision-variables related to inspection lead to updates in what information is available to the decision-maker at later time points, the formulated model deems the uncertainty in the program \textit{endogenous}, which necessitates scenario generation and reduction for tractability.
The program is implemented and solved for a smaller set of defects, showing how probabilistic forecasts can be connected to decision-making for infrastructure managers. Scaling up to a larger defect set increases the complexity of the scenario tree, and methods such as approximate dynamic programming can be used to handle the resulting curse of dimensionality \citep{powell2016a}. For validation, the solution method could be extended to a rolling-horizon setting, where updated predictive distributions are mapped to decisions at each stage, thereby giving a decision policy.\footnote{On policy-generation methods more broadly, we refer to \citet{powell2018a}.} In practice, this policy could be compared with the one currently in place at the infrastructure manager to validate the effectiveness of the decision support.
%Future research could focus on developing more sophisticated decision policies, for example by solving the proposed multistage program in a rolling-horizon setting. Alternative policy classes, such as those described in \citet{powell2018a}, may also be considered.

In summary, this study contributes both theoretically and practically: it advances Markov chain-based degradation modeling by introducing various model specifications and scoring-rule-consistent estimation, and it demonstrates how probabilistic forecasts can support infrastructure management. Beyond railways, the framework is broadly applicable to multicategorical forecasting problems in other domains where degradation processes evolve over ordinal state spaces and inspection intervals can be irregular.

